%=========================================
% PROOF OF COROLLARY 1
%=========================================

\section{Proof of Corollary 1}
\label{Proof of cor:error}

\subsection{Part (i)}

For the first quantile,
\begin{align*}
\frac{\mathcal{N}m_{1}}{\sqrt{T-R}}\sum_{t=R+1}^{T} \left( \widehat{S}_{\tau(t),R}^{p(1)}-S_{\tau(t)}^{p(1)}\right) &= \frac{\mathcal{N}m_{1}}{\sqrt{T-R}}\sum_{t=R+1}^{T} \left( \widehat{S}_{\tau(t),R}^{q(1)}-S_{\tau(t)}^{q(1)}\right) +o_{p}(1),
\end{align*}
where $\widehat{S}_{\tau(t),R}^{q(1)}$ and $S_{\tau(t)}^{q(1)}$ are the estimators of the first quantile based on the estimated and true sorting variable, i.e.
\begin{align*}
\widehat{S}_{\tau(t),R}^{q(1)}=\arg\min_{x}\frac{1}{N_{\tau(t)}}\sum_{i=1}^{N_{\tau(t)}}\left( \frac{1}{\mathcal{N}}-\mathds{1}\left\{ \widehat{S}_{i,\tau(t),R}\leq x\right\} \right)\left( \widehat{S}_{i,\tau(t),R}-x\right).
\end{align*}
\begin{align*}
S_{\tau(t)}^{q(1)}=\arg\min_{x}\frac{1}{N_{\tau(t)}}\sum_{i=1}^{N_{\tau(t)}}\left( \frac{1}{\mathcal{N}}-\mathds{1}\left\{ S_{i,\tau(t)}\leq x\right\} \right)\left( S_{i,\tau(t)}-x\right).
\end{align*}
Via an expansion of the pseudo score of $\widehat{S}_{\tau(t),R}^{q(1)}$ around the true quantile $S^{q(1)}$,
\begin{align*}
\widehat{S}_{\tau(t),R}^{q(1)}-S^{q(1)}
&=\frac{1}{N_{\tau(t)}}\sum_{i=1}^{N_{\tau(t)}}\left( \frac{1}{\mathcal{N}}-\mathds{1}\left\{ \widehat{S}_{i,\tau(t),R}\leq S^{q(1)}\right\} \right) \text{\hl{NEW:Not correct}} \\
&=\frac{1}{N_{\tau(t)}}\sum_{i=1}^{N_{\tau(t)}}\left( \frac{1}{\mathcal{N}}-\mathds{1}\left\{ S_{i,\tau(t)}\leq S^{q(1)}\right\} \right) \\
&+\frac{1}{N_{\tau(t)}}\sum_{i=1}^{N_{\tau(t)}}\left( \mathds{1}\left\{ \widehat{S}_{i,\tau(t),R}\leq S^{q(1)}\right\} -\mathds{1}\left\{ S_{i,\tau(t)}\leq S^{q(1)}\right\} \right) \text{\hl{NEW:Sign error}} \\
&=\frac{1}{N_{\tau(t)}}\sum_{i=1}^{N_{\tau(t)}}\left( \frac{1}{\mathcal{N}}-\mathds{1}\left\{ S_{i,\tau(t)}\leq S^{q(1)}\right\} \right) \\
&+\frac{1}{N_{\tau(t)}}\sum_{i=1}^{N_{\tau(t)}}\Bigg( \Big( \mathds{1}\left\{ S_{i,\tau(t)}\leq S^{q(1)} -\left( \widehat{S}_{i,\tau(t),R}-S_{i,\tau(t)}\right)\right\} -F\left( S^{q(1)}-\left( \widehat{S}_{i,\tau(t),R}-S_{i,\tau(t)}\right) \right) \Big) \\
&\hspace{1.7cm} -\Big( \mathds{1}\left\{ S_{i,\tau(t)}\leq S^{q(1)}\right\} - F\left( S^{q(1)}\right) \Big) \Bigg) \\
&+\frac{1}{N_{\tau(t)}}\sum_{i=1}^{N_{\tau(t)}}\left( F\left( S^{q(1)}-\left( \widehat{S}_{i,\tau(t),R}-S_{i,\tau(t)}\right) \right) -F\left( S^{q(1)}\right) \right) \\
&=I_{\tau(t),N_{\tau(t)}}+II_{\tau(t),N_{\tau(t)},R}+III_{\tau(t),N_{\tau(t)},R}.
\end{align*}
and note that $I_{\tau(t),N_{\tau(t)}}=S_{\tau(t)}^{q(1)}-S^{q(1)}$ \hl{Not sure why this holds} and that
\begin{equation*}
III_{\tau(t),N_{\tau(t)},R}=\frac{1}{N_{\tau(t)}}\sum_{i=1}^{N_{\tau(t)}}f\left( S^{q(1)}\right) \left( \widehat{S}_{i,\tau(t),R}-S_{i,\tau(t)}\right).
\end{equation*}
Moreover, given assumption \ref{ass:sups}, because of stochastic equicontinuity,
\begin{equation*}
\frac{\mathcal{N}}{\sqrt{T-R}}\sum_{t=R+1}^{T}III_{\tau(t),N_{\tau(t)},R}=o_{p}(I_{\tau(t),N_{\tau(t)},R} + II_{\tau(t),N_{\tau(t)},R}). \text{\hl{This doesn't look right}}
\end{equation*}

\hl{Where did $II_{\tau(t),N_{\tau(t)},R}$ disappear} Hence,
\begin{align*}
\frac{\mathcal{N}m_{1}}{\sqrt{T-R}}\sum_{t=R+1}^{T}\left( \widehat{S}_{\tau(t),R}^{q(1)}-S_{\tau(t)}^{q(1)}\right)
=\frac{\mathcal{N}m_{1}}{\sqrt{T-R}}\sum_{t=R+1}^{T}\frac{1}{N_{\tau(t)}}\sum_{i=1}^{N_{\tau(t)}}f\left( S^{q(1)}\right) \left( \widehat{S}_{i,\tau(t),R}-S_{i,\tau(t)}\right) +o_{p}(1).
\end{align*}
\hl{Here we have $\widehat{S}_{\tau(t),R}^{q(1)}-S_{\tau(t)}^{q(1)}$ but the LHS of the previous equation is $\widehat{S}_{\tau(t),R}^{q(1)}-S^{q(1)}$}

As before, we consider estimation error only for the first quantile. If $(T-R)/R\rightarrow 0,$
\begin{align*}
&\frac{\mathcal{N}}{\sqrt{T-R}}\sum_{t=R+1}^{T}m_{1}\left( \widehat{S}_{\tau(t),R}^{p(1)}-S_{\tau(t)}^{p(1)}\right) \\
&=\frac{\mathcal{N}}{\sqrt{T-R}}\sum_{t=R+1}^{T}m_{1}\left( \left( \widehat{S}_{\tau(t),R}^{p(1)}-S^{q(1)} \right) - \left( S_{\tau(t),R}^{p(1)}-S^{q(1)} \right) \right) \\
&=O_{p}\left( \sqrt{\frac{T-R}{R}}\right) =o_{p}(1),
\end{align*}
because,
\begin{align*}
\frac{\mathcal{N}}{\sqrt{T-R}}\sum_{t=R+1}^{T}\left( \widehat{S}_{\tau(t),R}^{p(1)}-S^{q(1)}\right)
=O_{p}\left( \sqrt{\frac{T-R}{R}}\right). \text{\hl{Check}}
\end{align*}
and, under assumption \refAssPart{ass:mixing, v},
\begin{align*}
\frac{\mathcal{N}}{\sqrt{T-R}}\sum_{t=R+1}^{T}\left( S_{\tau(t),R}^{p(1)}-S^{q(1)}\right)
=O_{p}\left( N^{-(1-\rho)/2} \sqrt{\frac{T-R}{R}}\right). \text{\hl{Check}}
\end{align*}
or, under assumption \ref{ass:strongdep},
\begin{align*}
\frac{\mathcal{N}}{\sqrt{T-R}}\sum_{t=R+1}^{T}\left( S_{\tau(t),R}^{p(1)}-S^{q(1)}\right)
=O_{p}\left( \sqrt{\frac{P}{T-R}}\right). \text{\hl{what is $P$?}}
\end{align*}
The statement in (\ref{No-PEE}) follows.

If instead $(T-R)/R\rightarrow \pi >0$, under assumption \refAssPart{ass:mixing, v}, \hl{Before we had $(\widehat{S}_{\tau(t),R}^{p(1)}-S^{q(1)}) - ( S_{\tau(t),R}^{p(1)}-S^{q(1)})$ now we have $(\widehat{S}_{\tau(t),R}^{p(1)}-S^{q(1)}_{\tau(t)} ) - ( S_{\tau(t),R}^{p(1)}-S^{q(1)}_{\tau(t)})$}
\begin{align*}
\frac{\mathcal{N}}{\sqrt{T-R}}\sum_{t=R+1}^{T}m_{1}\left( \left( \widehat{S}_{\tau(t),R}^{p(1)}-S^{q(1)}_{\tau(t)} \right) - \left( S_{\tau(t),R}^{p(1)}-S^{q(1)}_{\tau(t)} \right) \right) =o_{p}(1),
\end{align*}
and so that $\widehat{Z}_{T,N,T}-Z_{T,N}=o_p(1)$ and $Z_{T,N}$
converges in distribution to a normal variable with mean zero and variance equal to $V_{S^{q(1)}}$. Instead, under assumption \ref{ass:strongdep},
\begin{align*}
\frac{\mathcal{N}}{\sqrt{T-R}}\sum_{t=R+1}^{T}m_{1}\left( \left( \widehat{S}_{\tau(t),R}^{p(1)}-S^{q(1)}_{\tau(t)} \right) - \left( S_{\tau(t),R}^{p(1)}-S^{q(1)}_{\tau(t)} \right) \right)
\end{align*}
converges in distribution and it also contributes to the asymptotic variance.

\subsection{Part (ii)}

It follows immediately from Theorem~\ref{th:feasible}, part (ii).